\section{Uvod}

U savremenim informacionim tehnologijama, distribuirana tehnologija glavne knjige (DLT - engl. \textit{Distributed Ledger Technology}) predstavlja infrastrukturu za decentralizovano čuvanje i deljenje podataka među učesnicima mreže bez centralnog autoriteta. DLT omogućava da više računara (čvorova) poseduje identičnu evidenciju transakcija, gde se pre unošenja novih transakcija učesnici mreže dogovaraju o njihovom prihvatanju~\cite{blockchain and DLT}.

Prva i najpoznatija implementacija DLT-a je blokčejn (engl. \textit{blockchain}), koji je popularizovan kroz kriptovalutu \textit{Bitcoin} i mnoge druge projekte. U blokčejnu podaci se čuvaju u blokovima koji su linearno povezani i dodaju se kroz sekvencijalnu strukturu. Međutim, blokčejn ima svoje izazove, posebno u pogledu skalabilnosti, potrošnje električne energije i troškovima transakcija, što podstiče razvoj alternativnih arhitektura~\cite{blockchain and DLT}.

Jedna od ovih alternativa je usmereni aciklični graf (DAG - engl. \textit{Directed acyclic graph}), koji se koristi kao struktura podataka u određenim vrstama DLT sistema. U ovakvim sistemima, struktura usmerenog acikličnog grafa može u velikoj meri uticati na performanse mreže, kao i na način postizanja dogovora (konsenzusa) učesnika.

Cilj ovog rada je da prestavi koncept distribuirane glavne knjige bazirane na usmerenom acikličnom grafu, objasni njegovu osnovnu arhitekturu, prikaže razliku u odnosu na blokčejn, navede nekoliko realnih primera, kao i da diskutuje o osnovnim prednostima i izazovima koje ovakav pristup donosi.

%Struktura rada je organizovana tako da se najpre predstavi teorijska osnova usmerenog acikličnog grafa, zatim prikaže DLT modeliran kao DAG, objasni uspostavljanje konsenzusa u ovakvim sistemima i navede nekoliko realnih primera. Nakon toga će biti prikazane prednosti i ograničenja ovakvog pristupa, dok će u zaključku biti sumirani glavni nalazi rada.